\documentclass[12pt,a4paper]{article}

% ─── Packages ───────────────────────────────────────────────────────────────
\usepackage[utf8]{inputenc}
\usepackage[T1]{fontenc}
\usepackage{babel}
\usepackage{geometry}
\usepackage{graphicx}
\usepackage{float}
\usepackage{xcolor}
\usepackage{listings}
\usepackage{booktabs}
\usepackage{array}
\usepackage{longtable}
\usepackage{hyperref}
\usepackage{fancyhdr}
\usepackage{titlesec}
\usepackage{tcolorbox}
\usepackage{enumitem}
\usepackage{amsmath}
\usepackage{mdframed}
\usepackage{tabularx}

\geometry{
  top=2.5cm, bottom=2.5cm,
  left=2.5cm, right=2.5cm,
  headheight=15pt
}

% ─── Couleurs ────────────────────────────────────────────────────────────────
\definecolor{sambablue}{RGB}{0,82,156}
\definecolor{codebg}{RGB}{245,245,245}
\definecolor{codeframe}{RGB}{180,180,180}
\definecolor{alertred}{RGB}{180,0,0}
\definecolor{successgreen}{RGB}{0,130,0}
\definecolor{voipblue}{RGB}{0,90,160}
\definecolor{explainyellow}{RGB}{255,250,220}
\definecolor{explainframe}{RGB}{180,140,0}

% ─── Style des listings ──────────────────────────────────────────────────────
\lstdefinestyle{terminal}{
  backgroundcolor=\color{black},
  basicstyle=\ttfamily\small\color{white},
  breaklines=true,
  frame=single,
  rulecolor=\color{gray},
  xleftmargin=6pt,
  xrightmargin=6pt,
  aboveskip=8pt,
  belowskip=8pt,
}

\lstdefinestyle{config}{
  backgroundcolor=\color{codebg},
  basicstyle=\ttfamily\small,
  keywordstyle=\color{sambablue}\bfseries,
  commentstyle=\color{gray}\itshape,
  stringstyle=\color{alertred},
  breaklines=true,
  frame=single,
  rulecolor=\color{codeframe},
  xleftmargin=6pt,
  xrightmargin=6pt,
  aboveskip=8pt,
  belowskip=8pt,
  showstringspaces=false,
}

% ─── Boîtes colorées ─────────────────────────────────────────────────────────
\tcbuselibrary{skins,breakable}

\newtcolorbox{infobox}[1][]{
  colback=blue!5!white, colframe=sambablue,
  fonttitle=\bfseries, title=#1,
  breakable
}
\newtcolorbox{successbox}[1][]{
  colback=green!5!white, colframe=successgreen,
  fonttitle=\bfseries, title=#1,
  breakable
}
\newtcolorbox{warningbox}[1][]{
  colback=orange!8!white, colframe=orange!80!black,
  fonttitle=\bfseries, title=#1,
  breakable
}
\newtcolorbox{explainbox}[1][]{
  colback=explainyellow, colframe=explainframe,
  fonttitle=\bfseries\color{explainframe}, title=#1,
  breakable
}

% ─── En-têtes / pieds de page ────────────────────────────────────────────────
\pagestyle{fancy}
\fancyhf{}
\fancyhead[L]{\small Projet Final -- Serveur Samba PME}
\fancyhead[R]{\small \nouppercase{\leftmark}}
\fancyfoot[C]{\thepage}
\renewcommand{\headrulewidth}{0.4pt}

% ─── Titre des sections ──────────────────────────────────────────────────────
\titleformat{\section}
  {\Large\bfseries\color{sambablue}}{\thesection}{1em}{}[\titlerule]
\titleformat{\subsection}
  {\large\bfseries\color{sambablue!80!black}}{\thesubsection}{1em}{}
\titleformat{\subsubsection}
  {\normalsize\bfseries}{\thesubsubsection}{1em}{}

% ─── Liens hypertexte ────────────────────────────────────────────────────────
\hypersetup{
  colorlinks=true,
  linkcolor=sambablue,
  urlcolor=sambablue,
  citecolor=sambablue,
  pdftitle={Rapport Serveur Samba PME},
  pdfauthor={KPODONOU Kossigan Gael God-Love},
}

% ════════════════════════════════════════════════════════════════════════════
\begin{document}

% ─── PAGE DE TITRE ───────────────────────────────────────────────────────────
\begin{titlepage}
  \newgeometry{margin=2cm}

  % LOGOS en haut à gauche et à droite (positionnement absolu)
  \begin{tikzpicture}[remember picture, overlay]
    % Logo gauche
    \node[anchor=north west, xshift=1cm, yshift=-1cm] 
      at (current page.north west) {\includegraphics[width=5.5cm]{WhatsApp Image 2025-01-06 à 00.28.20_4ed6a7f1.jpg}};
    % Logo droit
    \node[anchor=north east, xshift=-1cm, yshift=-1cm] 
      at (current page.north east) {\includegraphics[width=5.5cm]{ecole_polytechnique.jpg}};
  \end{tikzpicture}

  \vspace*{4cm} 

  \vspace{2.5cm}

  \begin{tcolorbox}[
    colback=white, colframe=sambablue,
    boxrule=1.4pt, arc=5pt,
    left=22pt, right=22pt,
    top=16pt, bottom=16pt,
    width=\textwidth
  ]
    \centering
    {\LARGE\bfseries\color{voipblue} RAPPORT DE PROJET FINAL}\\[14pt]
    {\fontsize{19}{23}\selectfont\bfseries
      Mise en place d'un serveur\\[6pt]
      de fichiers Samba en PME}
  \end{tcolorbox}

  \vspace{3.5cm}

  \begin{tabular}{rl}
    \textbf{Auteur1}         & : KPODONOU Kossigan Gael God-Love \\[8pt]
    \textbf{Auteur2}         & : GNIGMA Bruno \\[8pt]
    \textbf{Professeur}     & : Dr. DJEKI \\[8pt]
    \textbf{Matière}        & : Généralités sur les Stockages de l'information (INF 1527)\\[8pt]
    \textbf{Année scolaire} & : 2025--2026 \\
  \end{tabular}

  \vfill
  
\end{titlepage}

% ─── TABLE DES MATIÈRES ──────────────────────────────────────────────────────
\tableofcontents
\newpage

% ════════════════════════════════════════════════════════════════════════════
\section{Introduction et objectifs}

Dans un contexte où la gestion centralisée des données est devenue indispensable pour toute
organisation, les PME font face à des défis croissants en matière de partage de fichiers, de
contrôle d'accès et de sécurité informatique. Ce projet a pour objet la conception et le
déploiement d'un serveur de fichiers \textbf{Samba} pour une PME fictive comptant quatre
services~: Direction, Comptabilité, Ressources Humaines et Technique.

\subsection*{Objectifs}

\begin{itemize}[leftmargin=2em]
  \item Centraliser les données de l'entreprise sur un serveur Linux sous Ubuntu Server~;
  \item Créer des espaces de stockage dédiés par service avec des droits d'accès stricts~;
  \item Sécuriser les communications en désactivant SMBv1 et en imposant SMB2/SMB3~;
  \item Protéger le serveur contre les attaques courantes (brute force, énumération anonyme)~;
  \item Mettre en place une stratégie de sauvegarde automatisée avec procédure de restauration testée~;
  \item Documenter l'intégralité de la démarche dans ce rapport.
\end{itemize}

% ════════════════════════════════════════════════════════════════════════════
\section{Recherche théorique}

\subsection{Qu'est-ce que Samba~?}

Samba est une implémentation libre et open-source du protocole \textbf{SMB/CIFS}
(\textit{Server Message Block / Common Internet File System}) pour les systèmes Unix/Linux.
Il permet à des machines Linux de partager des fichiers, des imprimantes et d'autres ressources
avec des clients Windows, macOS ou Linux sur un réseau local. Développé initialement par
Andrew Tridgell en 1992 par rétro-ingénierie du protocole SMB, Samba est aujourd'hui
maintenu par une large communauté open-source et constitue la solution de référence pour
l'interopérabilité Linux/Windows en entreprise.

Samba est composé de plusieurs démons~:
\begin{itemize}[leftmargin=2em]
  \item \texttt{smbd}~: gère les partages de fichiers, l'authentification des utilisateurs et
    les impressions réseau. C'est le démon principal que l'on interagit au quotidien~;
  \item \texttt{nmbd}~: gère la résolution de noms NetBIOS, permettant aux machines de
    se trouver par nom sur le réseau local~;
  \item \texttt{winbindd}~: permet l'intégration à un domaine Active Directory/LDAP,
    utile pour les grands réseaux d'entreprise.
\end{itemize}

\subsection{Le protocole SMB et ses versions}

Le protocole \textbf{SMB (Server Message Block)} est un protocole de couche application
permettant le partage de ressources (fichiers, imprimantes, ports série) entre machines d'un
réseau local. Initialement conçu par IBM et adopté par Microsoft, il a évolué de manière
significative au fil des années.

Dans l'ancien Windows NT~4, il était appelé \textbf{CIFS (Common Internet File System)}.
La version~2 de SMB est apparue dans Vista, Windows~7 et Windows~8. Actuellement
SMB est en version \textbf{3.1.1}, introduite dans Windows~10 et Windows Server~2016.

\begin{table}[H]
  \centering
  \begin{tabular}{lll}
    \toprule
    \textbf{Version} & \textbf{Systèmes concernés} & \textbf{Statut de sécurité} \\
    \midrule
    SMBv1 & Windows XP, Server 2003   & \textcolor{alertred}{\textbf{Obsolète -- à désactiver}} \\
    SMBv2 & Windows Vista, Server 2008 & Acceptable \\
    SMBv3 & Windows 8, Server 2012+   & \textcolor{successgreen}{\textbf{Recommandé}} \\
    \bottomrule
  \end{tabular}
  \caption{Versions du protocole SMB}
\end{table}

\begin{warningbox}[Danger SMBv1 -- Vulnérabilité critique]
SMBv1 est le protocole à l'origine des attaques \textit{WannaCry} (mai 2017) et
\textit{NotPetya} (juin 2017), qui ont causé des milliards de dollars de dommages dans le
monde entier, infectant des hôpitaux, des banques et des grandes entreprises. Sa
désactivation est une \textbf{exigence absolue de sécurité} dans tout environnement
professionnel.
\end{warningbox}

\subsection{Bonnes pratiques de sécurisation Samba}

\begin{enumerate}[leftmargin=2em]
  \item \textbf{Désactiver SMBv1}~: utiliser \texttt{min protocol = SMB2} dans la
    configuration globale pour forcer l'utilisation de protocoles modernes.
  \item \textbf{Authentification par utilisateur}~: \texttt{security = user} impose un
    login/mot de passe pour chaque connexion. C'est le mode opposé au mode ``share''
    (accès sans authentification) qui est dépassé.
  \item \textbf{Mots de passe robustes}~: chaque utilisateur doit avoir un mot de passe fort
    (majuscule, chiffre, caractère spécial, longueur $\geq$ 10 caractères) pour résister
    aux attaques par dictionnaire comme celle avec \texttt{rockyou.txt}.
  \item \textbf{Principe du moindre privilège}~: chaque groupe accède \textit{uniquement} à
    son dossier dédié. Un employé de la comptabilité ne doit pas pouvoir lire les fichiers RH.
  \item \textbf{Interdire la navigation anonyme}~: \texttt{restrict anonymous = 2} empêche
    un attaquant de lister les partages sans s'authentifier (énumération).
  \item \textbf{Journalisation}~: activer les logs (niveau 2) pour tracer toutes les
    connexions et tentatives d'accès. Indispensable en cas d'incident.
  \item \textbf{Protection contre le brute force}~: déployer \texttt{fail2ban} pour bannir
    automatiquement les IP suspectes après un nombre d'échecs configurable.
  \item \textbf{Sauvegarde régulière}~: mettre en place des sauvegardes complètes et
    incrémentielles automatisées, avec tests de restauration réguliers.
\end{enumerate}

\subsection{Permissions Linux et Samba}

La sécurité d'un partage Samba repose sur \textbf{deux couches de permissions} qui
s'appliquent simultanément et de manière indépendante~:

\begin{itemize}[leftmargin=2em]
  \item \textbf{Permissions Linux (POSIX)}~: définies via \texttt{chmod} et \texttt{chown}.
    Ce sont les permissions du système de fichiers sous-jacent, appliquées par le noyau Linux
    indépendamment de Samba.
  \item \textbf{Permissions Samba}~: définies dans \texttt{smb.conf} via \texttt{valid users},
    \texttt{read only}, \texttt{create mask}, etc.
\end{itemize}

\begin{infobox}[Règle des permissions effectives]
La permission effective est le \textbf{plus restrictif} des deux niveaux. Un utilisateur qui
a accès via Samba (\texttt{valid users}) mais dont le compte Linux n'a pas les droits sur
le dossier se verra quand même refusé. Il est donc indispensable de configurer
correctement \textit{les deux} niveaux.
\end{infobox}

% ════════════════════════════════════════════════════════════════════════════
\section{Architecture et environnement}

\subsection{Topologie du laboratoire}

L'environnement de test a été mis en place sur VirtualBox installé sur un hôte physique
fonctionnant sous Ubuntu Desktop. Trois machines constituent l'infrastructure~:

\begin{table}[H]
  \centering
  \begin{tabular}{llll}
    \toprule
    \textbf{Rôle} & \textbf{Système} & \textbf{IP} & \textbf{Interface} \\
    \midrule
    Serveur Samba & Ubuntu Server 24.04 & 192.168.1.118 (Bridge)    & enp0s3 \\
                  &                     & 192.168.56.20 (Host-Only) & enp0s8 \\
    Client Linux  & Ubuntu Desktop (hôte) & 192.168.56.1            & vboxnet0 \\
    Attaquant     & Kali Linux 2025.4   & 192.168.56.40             & eth1 \\
    \bottomrule
  \end{tabular}
  \caption{Inventaire des machines du laboratoire}
\end{table}

\subsection{Choix de la configuration réseau}

Deux cartes réseau ont été configurées sur la VM serveur~:

\begin{itemize}[leftmargin=2em]
  \item \textbf{Carte 1 (Bridge -- enp0s3)}~: place la VM sur le même réseau physique que
    la machine hôte. Permet l'accès à Internet pour l'installation des paquets
    (\texttt{apt install}). Elle reçoit son IP via DHCP depuis la box domestique.
  \item \textbf{Carte 2 (Host-Only -- enp0s8)}~: crée un réseau privé
    \texttt{192.168.56.0/24} entre la VM, l'hôte et Kali. C'est sur cette interface que
    transitent les connexions Samba et les tests de sécurité.
\end{itemize}

\begin{infobox}[Pourquoi Host-Only pour les tests Samba~?]
Le réseau Host-Only isole le trafic de test du réseau physique, ce qui est une bonne
pratique en laboratoire. Cela évite d'exposer le serveur (potentiellement non complètement
sécurisé en phase de test) aux autres machines du réseau réel, et empêche tout risque de
perturbation du réseau domestique par les attaques générées depuis Kali.
\end{infobox}

\subsubsection{Configuration réseau du serveur (netplan)}

Ubuntu Server utilise \textbf{netplan} pour gérer les interfaces réseau. La configuration
est déclarative (format YAML) et remplace les anciens fichiers \texttt{/etc/network/interfaces}.

\begin{lstlisting}[style=config, caption={/etc/netplan/00-installer-config.yaml}]
network:
  version: 2
  ethernets:
    enp0s3:
      dhcp4: true        # IP automatique via DHCP (acces internet)
    enp0s8:
      dhcp4: no          # Pas de DHCP : IP statique
      addresses:
        - 192.168.56.20/24  # IP fixe sur le reseau Host-Only
\end{lstlisting}

\begin{explainbox}[Explication~: \texttt{netplan apply}]
Après toute modification du fichier YAML, la commande \texttt{sudo netplan apply}
recharge la configuration réseau sans redémarrer la machine. Si l'interface reste en état
\texttt{NO-CARRIER}, vérifier que la case \textit{Câble connecté} est cochée dans les
paramètres VirtualBox de l'adaptateur réseau.
\end{explainbox}

\subsection{Organisation des services et utilisateurs}

Le projet impose la création d'au moins 8 utilisateurs répartis en 4 groupes correspondant
aux services de la PME.

\begin{table}[H]
  \centering
  \begin{tabular}{lll}
    \toprule
    \textbf{Groupe Linux} & \textbf{Utilisateurs}      & \textbf{Dossier partage} \\
    \midrule
    \texttt{direction}    & direction1, direction2     & \texttt{/srv/samba/direction} \\
    \texttt{comptabilite} & compta1, compta2            & \texttt{/srv/samba/comptabilite} \\
    \texttt{rh}           & rh1, rh2                    & \texttt{/srv/samba/rh} \\
    \texttt{technique}    & tech1, tech2                & \texttt{/srv/samba/technique} \\
    \textit{(tous)}       & tous les groupes            & \texttt{/srv/samba/commun} \\
    \bottomrule
  \end{tabular}
  \caption{Organisation des groupes, utilisateurs et partages}
\end{table}

% ════════════════════════════════════════════════════════════════════════════
\section{Installation et configuration}

\subsection{Installation de Samba}

Après mise à jour du système, Samba est installé via le gestionnaire de paquets~:

\begin{lstlisting}[style=terminal]
sudo apt update && sudo apt install samba -y
samba --version
sudo systemctl status smbd
\end{lstlisting}

\begin{figure}[H]
  \centering
  \includegraphics[width=1\textwidth]{samba_actif.png}
  \caption{Samba 4.19.5 installé et service \texttt{smbd} actif -- statut \texttt{active (running)}}
\end{figure}

\begin{explainbox}[Explication~: \texttt{sudo apt update \&\& sudo apt install samba -y}]
\begin{itemize}[leftmargin=1.5em,topsep=2pt]
  \item \texttt{sudo}~: exécute la commande avec les privilèges administrateur (root).
  \item \texttt{apt update}~: met à jour la liste des paquets disponibles depuis les dépôts
    officiels. Ne met PAS à jour les paquets installés -- seulement la liste.
  \item \texttt{\&\&}~: enchaîne la commande suivante \textit{uniquement} si la précédente
    a réussi (code de retour 0).
  \item \texttt{apt install samba -y}~: installe le paquet \texttt{samba} et ses dépendances.
    Le \texttt{-y} répond automatiquement ``oui'' aux questions de confirmation.
\end{itemize}
\end{explainbox}

\subsection{Création de la structure des dossiers}

\begin{lstlisting}[style=terminal]
sudo mkdir -p /srv/samba/direction
sudo mkdir -p /srv/samba/comptabilite
sudo mkdir -p /srv/samba/rh
sudo mkdir -p /srv/samba/technique
sudo mkdir -p /srv/samba/commun
\end{lstlisting}

\begin{explainbox}[Explication~: \texttt{mkdir -p}]
\begin{itemize}[leftmargin=1.5em,topsep=2pt]
  \item \texttt{mkdir}~: crée un répertoire (\textit{make directory}).
  \item \texttt{-p}~: crée les répertoires parents manquants sans erreur. Ici, si
    \texttt{/srv/samba/} n'existe pas encore, il sera créé automatiquement avant
    \texttt{direction/}. Sans \texttt{-p}, la commande échouerait si le parent n'existe pas.
  \item \texttt{/srv/samba/}~: chemin standardisé FHS (Filesystem Hierarchy Standard)
    pour les données de services. \texttt{/srv} est prévu pour les données servies par le
    système.
\end{itemize}
\end{explainbox}

\subsection{Création des groupes et utilisateurs}

\subsubsection{Création des groupes Linux}

\begin{lstlisting}[style=terminal]
sudo groupadd direction
sudo groupadd comptabilite
sudo groupadd rh
sudo groupadd technique
\end{lstlisting}

\begin{explainbox}[Explication~: \texttt{groupadd}]
\begin{itemize}[leftmargin=1.5em,topsep=2pt]
  \item \texttt{groupadd}~: crée un nouveau groupe dans le fichier \texttt{/etc/group}.
  \item Chaque groupe se voit attribuer un \textbf{GID} (Group IDentifier) unique,
    automatiquement incrémenté à partir de 1000 pour les groupes système non-root.
  \item Les groupes Linux sont le mécanisme fondamental de contrôle d'accès. Un fichier
    appartenant au groupe \texttt{direction} n'est accessible qu'aux membres de ce groupe
    (si les permissions le définissent ainsi).
  \item Vérification~: \texttt{cat /etc/group | grep direction} pour voir le GID attribué.
\end{itemize}
\end{explainbox}

\subsubsection{Création des utilisateurs Linux et Samba}

Chaque utilisateur est créé avec son groupe principal, un mot de passe Linux
(\texttt{passwd}), puis enregistré dans la base Samba (\texttt{smbpasswd})~:

\begin{lstlisting}[style=terminal]
sudo useradd -m -G direction direction1
sudo useradd -m -G direction direction2
sudo useradd -m -G comptabilite compta1
sudo useradd -m -G comptabilite compta2
sudo useradd -m -G rh rh1
sudo useradd -m -G rh rh2
sudo useradd -m -G technique tech1
sudo useradd -m -G technique tech2

# Mot de passe Linux (exemple pour direction1)
sudo passwd direction1   # ex: Direction1/+2026

# Enregistrement dans la base Samba
sudo smbpasswd -a direction1
# Repeter pour chaque utilisateur
\end{lstlisting}

\begin{explainbox}[Explication~: \texttt{useradd -m -G <groupe> <nom>}]
\begin{itemize}[leftmargin=1.5em,topsep=2pt]
  \item \texttt{useradd}~: crée un nouvel utilisateur dans \texttt{/etc/passwd} et
    \texttt{/etc/shadow}.
  \item \texttt{-m}~: crée le répertoire personnel \texttt{/home/direction1/} de l'utilisateur.
    Sans \texttt{-m}, aucun home n'est créé, ce qui peut causer des problèmes de connexion.
  \item \texttt{-G direction}~: ajoute l'utilisateur au groupe supplémentaire \texttt{direction}.
    Distinction importante~: le groupe \textit{principal} (défini par \texttt{-g}) est celui
    par défaut lors de la création de fichiers~; les groupes \textit{supplémentaires}
    (définis par \texttt{-G}) donnent des droits d'accès additionnels.
  \item \texttt{sudo passwd direction1}~: définit le mot de passe du compte Linux. Ce mot
    de passe est stocké (haché) dans \texttt{/etc/shadow}.
\end{itemize}
\end{explainbox}

\begin{infobox}[Deux bases de mots de passe indépendantes]
Samba maintient sa \textbf{propre base de mots de passe} dans
\texttt{/var/lib/samba/private/passdb.tdb}, totalement indépendante du système Linux
(\texttt{/etc/shadow}). La commande \texttt{smbpasswd -a <user>} enregistre l'utilisateur
dans cette base Samba. Un utilisateur peut avoir un mot de passe Linux différent de son
mot de passe Samba. Pour simplifier l'administration en laboratoire, les mêmes mots de
passe ont été utilisés dans les deux bases.
\end{infobox}

\subsection{Application des permissions sur les dossiers}

\begin{lstlisting}[style=terminal]
# Etape 1 : Changer le proprietaire (groupe) de chaque dossier
sudo chown root:direction    /srv/samba/direction
sudo chown root:comptabilite /srv/samba/comptabilite
sudo chown root:rh           /srv/samba/rh
sudo chown root:technique    /srv/samba/technique
sudo chown root:root         /srv/samba/commun

# Etape 2 : Appliquer les permissions
sudo chmod 2770 /srv/samba/direction
sudo chmod 2770 /srv/samba/comptabilite
sudo chmod 2770 /srv/samba/rh
sudo chmod 2770 /srv/samba/technique
sudo chmod 2777 /srv/samba/commun
\end{lstlisting}

\begin{explainbox}[Explication détaillée~: \texttt{chown root:direction /srv/samba/direction}]
\begin{itemize}[leftmargin=1.5em,topsep=2pt]
  \item \texttt{chown}~: change le propriétaire (\textit{change owner}) d'un fichier ou
    répertoire.
  \item Syntaxe~: \texttt{chown <utilisateur>:<groupe> <cible>}
  \item \texttt{root:direction}~: le propriétaire est \texttt{root} (seul l'admin peut
    modifier les permissions globales), et le groupe propriétaire est \texttt{direction}
    (les membres de ce groupe auront les droits définis pour le groupe).
  \item Cela signifie~: \textit{ce dossier appartient à root mais est administré par le
    groupe direction}.
\end{itemize}
\end{explainbox}

\begin{explainbox}[Explication détaillée~: \texttt{chmod 2770} et \texttt{chmod 2777}]
Les permissions Linux s'expriment en octal sur 4 chiffres~:

\begin{center}
\begin{tabular}{|c|c|c|c|}
  \hline
  \textbf{Bit spécial} & \textbf{Propriétaire} & \textbf{Groupe} & \textbf{Autres} \\
  \hline
  2 (setgid) & 7 (rwx) & 7 (rwx) & 0 (---) \\
  \hline
\end{tabular}
\end{center}

\vspace{6pt}
\begin{itemize}[leftmargin=1.5em,topsep=2pt]
  \item \textbf{Valeur 7 = rwx}~: r (read=4) + w (write=2) + x (execute=1) = 7.
    Accès complet~: lecture, écriture, navigation.
  \item \textbf{Valeur 0 = ---}~: aucun droit. Les \textit{autres} (utilisateurs hors
    propriétaire et hors groupe) ne peuvent rien faire.
  \item \textbf{Bit setgid (2 en tête)}~: quand un utilisateur crée un fichier dans ce
    dossier, le fichier hérite automatiquement du \textbf{groupe du dossier}
    (ex.~\texttt{direction}) et non du groupe principal de l'utilisateur. Cela garantit
    que tous les fichiers créés restent accessibles à tout le groupe.
  \item \textbf{chmod 2777} (dossier commun)~: tout le monde (propriétaire, groupe, autres)
    a tous les droits. Avec le setgid, les nouveaux fichiers appartiennent au groupe du
    dossier. Cela permet à tous les employés (tous groupes confondus) de partager des
    fichiers dans l'espace commun.
  \item \textbf{Résumé visuel}~:
    \begin{itemize}[leftmargin=1.5em]
      \item \texttt{2770} $\rightarrow$ \texttt{drwxrws---}~: accès total groupe, zéro pour les autres
      \item \texttt{2777} $\rightarrow$ \texttt{drwxrwsrwx}~: accès total pour tous + setgid
    \end{itemize}
\end{itemize}
\end{explainbox}

\begin{figure}[H]
  \centering
  \includegraphics[width=0.88\textwidth]{permissions_dossiers.png}
  \caption{Structure \texttt{/srv/samba/} -- permissions correctement appliquées.
    On distingue clairement les groupes propriétaires (\texttt{direction},
    \texttt{comptabilite}, \texttt{rh}, \texttt{technique}) et les permissions
    \texttt{drwxrws---} (2770) et \texttt{drwxrwsrwx} (2777) pour \texttt{commun}.}
\end{figure}

\subsection{Configuration du fichier smb.conf}

Le fichier de configuration principal est sauvegardé avant modification (bonne pratique
indispensable)~:

\begin{lstlisting}[style=terminal]
sudo cp /etc/samba/smb.conf /etc/samba/smb.conf.bak
sudo nano /etc/samba/smb.conf
\end{lstlisting}

\begin{lstlisting}[style=config, caption={/etc/samba/smb.conf -- configuration complète}]
[global]
   workgroup = WORKGROUP
   server string = Serveur Samba PME
   server role = standalone server

   # === SECURITE ===
   security = user          # Authentification par compte utilisateur
   min protocol = SMB2      # Interdit SMBv1 (WannaCry, NotPetya)
   max protocol = SMB3      # Limite au protocole moderne

   # Interdit la navigation anonyme (enumeration des partages)
   restrict anonymous = 2

   # === LOGS ===
   log file = /var/log/samba/log.%m  # Un fichier par machine cliente
   max log size = 1000               # Taille max en Ko
   logging = file
   log level = 2                     # Niveau detaille (0=minimal, 10=max)

[direction]
   path = /srv/samba/direction
   valid users = @direction   # @ = groupe Linux
   read only = no             # Ecriture autorisee
   browseable = yes
   create mask = 0660         # Nouveaux fichiers : rw-rw----
   directory mask = 2770      # Nouveaux dossiers : rwxrws---

[comptabilite]
   path = /srv/samba/comptabilite
   valid users = @comptabilite
   read only = no
   browseable = yes
   create mask = 0660
   directory mask = 2770

[rh]
   path = /srv/samba/rh
   valid users = @rh
   read only = no
   browseable = yes
   create mask = 0660
   directory mask = 2770

[technique]
   path = /srv/samba/technique
   valid users = @technique
   read only = no
   browseable = yes
   create mask = 0660
   directory mask = 2770

[commun]
   path = /srv/samba/commun
   valid users = @direction @comptabilite @rh @technique
   read only = no
   browseable = yes
   create mask = 0666   # Fichiers lisibles/modifiables par tous
   directory mask = 2777
\end{lstlisting}

\begin{explainbox}[Explication des directives clés de smb.conf]
\begin{itemize}[leftmargin=1.5em,topsep=2pt]
  \item \textbf{\texttt{security = user}}~: chaque connexion doit s'authentifier avec
    un nom d'utilisateur et un mot de passe Samba. Mode opposé~: \texttt{security = share}
    (accès sans mot de passe, dépassé et dangereux).
  \item \textbf{\texttt{min protocol = SMB2}}~: force l'utilisation de SMB version 2 minimum.
    Toute tentative de connexion en SMBv1 est rejetée.
  \item \textbf{\texttt{restrict anonymous = 2}}~: niveau maximum de restriction anonyme.
    Valeur 0 = pas de restriction, valeur 1 = restriction partielle, valeur 2 = aucun accès
    anonyme du tout (énumération impossible).
  \item \textbf{\texttt{valid users = @direction}}~: le \texttt{@} indique un groupe
    Linux. Seuls les membres du groupe \texttt{direction} peuvent accéder à ce partage.
    Sans \texttt{@}, ce serait un utilisateur individuel.
  \item \textbf{\texttt{create mask = 0660}}~: masque de permissions pour les nouveaux
    fichiers. \texttt{0660} = rw-rw---- (propriétaire et groupe peuvent lire/écrire,
    les autres n'ont rien). Équivalent de \texttt{chmod 660} sur chaque nouveau fichier.
  \item \textbf{\texttt{log.~\%m}}~: \texttt{\%m} est une variable Samba remplacée par le
    nom NetBIOS ou l'adresse IP du client. Cela crée un fichier de log par machine
    connectée, facilitant les investigations.
\end{itemize}
\end{explainbox}

\subsubsection{Validation de la configuration}

\begin{lstlisting}[style=terminal]
testparm               # Verifie smb.conf sans redemarrer Samba
sudo systemctl restart smbd  # Prend en compte les modifications
\end{lstlisting}

\begin{figure}[H]
  \centering
  \includegraphics[width=0.88\textwidth]{vrai_test_fichier.png}
  \caption{\texttt{testparm} -- configuration valide. Le message
    \texttt{Loaded services file OK} confirme l'absence d'erreur de syntaxe.
    L'avertissement \texttt{Weak crypto is allowed by GnuTLS} est informatif
    (compatibilité NTLMv2) et n'est pas bloquant.}
\end{figure}

\begin{figure}[H]
  \centering
  \includegraphics[width=0.88\textwidth]{vrai_confirmation_reussie.png}
  \caption{Résultat de \texttt{testparm} affichant le dump complet de la configuration --
    on reconnaît les sections \texttt{[direction]}, \texttt{[comptabilite]},
    \texttt{[rh]}, \texttt{[technique]}, \texttt{[commun]} avec leurs paramètres.
    Service \texttt{smbd} ensuite redémarré avec succès.}
\end{figure}

% ════════════════════════════════════════════════════════════════════════════
\section{Tests de validation des accès}

\subsection{Test de connexion réussie -- partage \texttt{direction}}

Depuis le PC hôte (client Linux), on utilise \texttt{smbclient} pour se connecter~:

\begin{lstlisting}[style=terminal]
smbclient //192.168.56.20/direction -U direction1
\end{lstlisting}

\begin{figure}[H]
  \centering
  \includegraphics[width=1\textwidth]{test_connexion_samba.png}
  \caption{Première tentative avec SMB1 -- message \texttt{SMB1 disabled -- no workgroup
    available} confirme que le serveur rejette correctement les connexions SMBv1.
    C'est le comportement attendu de notre configuration \texttt{min protocol = SMB2}.}

  \vspace{0.4cm}
  \includegraphics[width=1\textwidth]{test_connexion_securite_reussi.png}
  \caption{Connexion SMB2 réussie par \texttt{direction1}~: création du dossier
    \texttt{test\_direction} visible via \texttt{ls}, suivi du test d'isolation
    (\texttt{smbclient .../rh}) qui retourne \texttt{NT\_STATUS\_ACCESS\_DENIED}.
    Les deux tests sont visibles dans la même capture.}
  \label{fig:connexion_reussie}
\end{figure}

\subsection{Test de refus d'accès -- isolation des partages}

On tente d'accéder au partage \texttt{rh} avec l'utilisateur \texttt{direction1} qui
n'appartient pas au groupe \texttt{rh}~:

\begin{lstlisting}[style=terminal]
smbclient //192.168.56.20/rh -U direction1
\end{lstlisting}

\begin{successbox}[Résultat attendu -- isolation confirmée]
\texttt{tree connect failed: NT\_STATUS\_ACCESS\_DENIED}\\[4pt]
L'accès est bien refusé. Les partages sont correctement isolés entre les groupes.
Ce résultat est visible dans la figure~\ref{fig:connexion_reussie} ci-dessus (dernière ligne
de la capture).
\end{successbox}

\subsection{Test du partage commun -- accès multi-groupes}

\textbf{Étape 1}~: \texttt{direction1} dépose un fichier dans le partage commun~:

\begin{lstlisting}[style=terminal]
# Creation du fichier sur la machine hote
echo "Fichier de test PME" > /tmp/fichier_test.txt

# Connexion au partage commun et upload
smbclient //192.168.56.20/commun -U direction1
smb: \> put /tmp/fichier_test.txt fichier_test.txt
smb: \> ls
\end{lstlisting}

\begin{figure}[H]
  \centering
  \includegraphics[width=0.88\textwidth]{test_validation-1.png}
  \caption{Upload du fichier \texttt{fichier\_test.txt} dans \texttt{/commun} par
    \texttt{direction1}. Le listing \texttt{ls} confirme sa présence (taille 20 octets,
    date de création visible).}
\end{figure}

\textbf{Étape 2}~: \texttt{compta1} (groupe différent) peut lire et télécharger ce fichier~:

\begin{figure}[H]
  \centering
  \includegraphics[width=0.88\textwidth]{fichier_test_accessible.png}
  \caption{\texttt{compta1} se connecte à \texttt{/commun}, liste les fichiers et
    télécharge \texttt{fichier\_test.txt} via \texttt{get}. Le fichier déposé par un membre
    d'un groupe différent est bien accessible. Le partage commun est fonctionnel.}
\end{figure}

\begin{table}[H]
  \centering
  \begin{tabular}{lccccc}
    \toprule
    \textbf{Utilisateur} & \textbf{/direction} & \textbf{/comptabilite} & \textbf{/rh} & \textbf{/technique} & \textbf{/commun} \\
    \midrule
    direction1 & \textcolor{successgreen}{Oui} & \textcolor{alertred}{Non} & \textcolor{alertred}{Non} & \textcolor{alertred}{Non} & \textcolor{successgreen}{Oui} \\
    compta1    & \textcolor{alertred}{Non} & \textcolor{successgreen}{Oui} & \textcolor{alertred}{Non} & \textcolor{alertred}{Non} & \textcolor{successgreen}{Oui} \\
    rh1        & \textcolor{alertred}{Non} & \textcolor{alertred}{Non} & \textcolor{successgreen}{Oui} & \textcolor{alertred}{Non} & \textcolor{successgreen}{Oui} \\
    tech1      & \textcolor{alertred}{Non} & \textcolor{alertred}{Non} & \textcolor{alertred}{Non} & \textcolor{successgreen}{Oui} & \textcolor{successgreen}{Oui} \\
    \bottomrule
  \end{tabular}
  \caption{Matrice des accès -- Oui = autorisé, Non = refusé (NT\_STATUS\_ACCESS\_DENIED)}
\end{table}

% ════════════════════════════════════════════════════════════════════════════
\section{Journalisation}

\subsection{Configuration et vérification}

La journalisation est activée dans \texttt{smb.conf} avec le niveau 2, ce qui enregistre
les authentifications, les connexions et les accès aux partages. Les logs sont générés dans
\texttt{/var/log/samba/} avec un fichier par machine cliente (\texttt{log.\%m}).

\begin{figure}[H]
  \centering
  \includegraphics[width=0.88\textwidth]{verification_logs.png}
  \vspace{0.4cm}
  \includegraphics[width=0.88\textwidth]{voir_log-1.png}
  \caption{Contenu de \texttt{/var/log/samba/} après les tests~: on distingue un fichier
    par machine ayant communiqué avec le serveur (\texttt{log.192.168.56.1} pour le
    client hôte, \texttt{log.192.168.56.30} et \texttt{log.192.168.56.40} pour les
    attaquants Kali). Le fichier \texttt{log.adrien-rc530-rc730} correspond au nom
    NetBIOS de la machine hôte résolue par Samba.}
\end{figure}

\subsection{Contenu des logs -- traçabilité des connexions}

Le fichier \texttt{log.192.168.56.1} (machine hôte) contient les traces d'authentification~:

\begin{figure}[H]
  \centering
  \includegraphics[width=0.88\textwidth]{connexion_client_hote.png}
  \caption{Extrait du fichier \texttt{log.192.168.56.1}~: chaque entrée indique
    l'horodatage précis, le résultat de l'authentification (\texttt{succeeded}),
    le protocole SMB2/NTLMv2, et les partages chargés lors de la session
    (\texttt{direction}, \texttt{comptabilite}, \texttt{rh}, \texttt{technique}, \texttt{commun}).}
\end{figure}

Chaque entrée de log contient~:
\begin{itemize}[leftmargin=2em]
  \item L'horodatage précis de la connexion (format \texttt{[AAAA/MM/JJ HH:MM:SS]})~;
  \item L'utilisateur ayant tenté la connexion~;
  \item Le résultat (\texttt{succeeded} ou \texttt{FAILED with error NT\_STATUS\_WRONG\_PASSWORD})~;
  \item Le protocole utilisé (SMB2/NTLMv2)~;
  \item L'adresse IP et le port de la machine cliente~;
  \item Les partages accédés durant la session.
\end{itemize}

\begin{infobox}[Utilité forensique des logs]
En cas d'incident de sécurité, ces logs permettent de retracer précisément~: qui s'est
connecté, quand, depuis quelle machine, sur quel partage, et si la tentative a réussi ou
échoué. C'est un élément indispensable pour toute investigation post-incident.
\end{infobox}

% ════════════════════════════════════════════════════════════════════════════
\section{Sécurité -- Protection contre les attaques}

\subsection{Désactivation de SMBv1}

SMBv1 est désactivé via la directive \texttt{min protocol = SMB2} dans la section
\texttt{[global]} de \texttt{smb.conf}. Cette mesure est la première et la plus fondamentale
des protections. Elle a été validée lors des tests d'attaque depuis Kali (section~7.3).

\subsection{Attaque 1 -- Énumération anonyme}

\subsubsection{Démonstration de la vulnérabilité (avant correction)}

Depuis Kali Linux, un attaquant tente de lister les partages \textit{sans s'authentifier}~:

\begin{lstlisting}[style=terminal]
smbclient -L //192.168.56.20 -N
# -L : List shares (lister les partages)
# -N : No password (connexion anonyme sans mot de passe)
\end{lstlisting}

\begin{figure}[H]
  \centering
  \includegraphics[width=0.88\textwidth]{partage_sans_authenetification.png}
  \caption{AVANT correction~: \texttt{Anonymous login successful} -- la liste complète
    des partages (\texttt{direction}, \texttt{comptabilite}, \texttt{rh},
    \texttt{technique}, \texttt{commun}) est visible sans aucun mot de passe.
    Un attaquant connaît déjà l'architecture interne de l'entreprise.}
\end{figure}

\begin{warningbox}[Vulnérabilité détectée -- Fuite d'information critique]
L'accès anonyme permettait à n'importe quelle machine du réseau de découvrir
l'existence et les noms de tous les partages sans la moindre authentification. Dans un
contexte réel, un attaquant obtient ainsi une carte de l'organisation interne de
l'entreprise, ce qui facilite les attaques ciblées (phishing, ingénierie sociale).
\end{warningbox}

\subsubsection{Correction -- restriction de l'accès anonyme}

Ajout de la directive suivante dans la section \texttt{[global]} de \texttt{smb.conf}~:

\begin{lstlisting}[style=config]
[global]
   ...
   restrict anonymous = 2
\end{lstlisting}

Puis redémarrage du service~: \texttt{sudo systemctl restart smbd}

\subsubsection{Validation de la correction}

\begin{figure}[H]
  \centering
  \includegraphics[width=0.88\textwidth]{anonymous_refuse.png}
  \caption{APRÈS correction~: \texttt{Anonymous login successful} s'affiche toujours
    (ouverture de session initiale autorisée) mais la tentative de lister les partages
    retourne immédiatement \texttt{tree connect failed: NT\_STATUS\_ACCESS\_DENIED}.
    L'énumération est bloquée.}
\end{figure}

\begin{successbox}[Protection validée -- énumération impossible]
Malgré l'ouverture de session anonyme initiale (comportement SMB normal), la tentative
de lister les partages est bloquée par \texttt{NT\_STATUS\_ACCESS\_DENIED}. Un
attaquant ne peut plus découvrir l'architecture des partages sans s'authentifier avec
un compte valide.
\end{successbox}

\subsection{Attaque 2 -- Brute force par dictionnaire (Hydra / CrackMapExec)}

\subsubsection{Tentative avec Hydra}

L'outil Hydra tente une attaque par brute force sur le service SMB avec le dictionnaire
\texttt{rockyou.txt} (14 millions de mots de passe communs)~:

\begin{lstlisting}[style=terminal]
hydra -l direction1 -P /usr/share/wordlists/rockyou.txt 192.168.56.20 smb
# -l : login cible (un seul utilisateur)
# -P : fichier de mots de passe (dictionnaire)
\end{lstlisting}

\begin{figure}[H]
  \centering
  \includegraphics[width=0.75\textwidth]{echec_hydra_du_a_smb1.png}
  \caption{Échec d'Hydra~: \texttt{[ERROR] target does not support SMBv1}. Hydra
    utilise SMBv1 par défaut pour les attaques SMB~; notre configuration
    \texttt{min protocol = SMB2} le bloque immédiatement.}

  \vspace{0.4cm}
  \includegraphics[width=1\textwidth]{hydra_echoue_encore.png}
  \caption{Même en forçant SMB2 avec le flag \texttt{-m smb2}, Hydra échoue encore~:
    \texttt{[ERROR] target smb://192.168.56.20:445/ does not support SMBv1}.
    L'outil ne supporte pas correctement le mode SMB2 natif pour le brute force.}
\end{figure}

\subsubsection{Attaque avec CrackMapExec (SMB2/NTLMv2)}

CrackMapExec est un outil plus moderne supportant nativement SMB2~:

\begin{lstlisting}[style=terminal]
crackmapexec smb 192.168.56.20 -u direction1 \
    -p /usr/share/wordlists/rockyou.txt
\end{lstlisting}

\begin{figure}[H]
  \centering
  \includegraphics[width=\textwidth]{kali_attaque_ubuntu_repertorie.png}
  \caption{Attaque brute force CrackMapExec en cours~: à gauche les logs Samba en
    temps réel (centaines de \texttt{STATUS\_LOGON\_FAILURE}), à droite CrackMapExec
    essayant chaque mot de passe du dictionnaire. L'IP attaquante
    \texttt{ipv4:192.168.56.30} est clairement identifiable dans les logs. Aucun mot de
    passe n'est trouvé car il respecte les critères de robustesse.}
\end{figure}

\begin{infobox}[Pourquoi l'attaque échoue~?]
Le mot de passe \texttt{Direction1/+2026} ne figure pas dans \texttt{rockyou.txt} car
il contient une majuscule, des chiffres et des caractères spéciaux. Les dictionnaires
de brute force contiennent des mots courants, des noms propres et des combinaisons
simples -- pas des mots de passe complexes et uniques. C'est pourquoi la politique de
mots de passe robustes est fondamentale.
\end{infobox}

\subsection{Fail2ban -- blocage automatique des attaquants}

\subsubsection{Installation et configuration}

\begin{lstlisting}[style=terminal]
sudo apt install fail2ban -y
\end{lstlisting}

\textbf{Fichier} \texttt{/etc/fail2ban/jail.local}~:

\begin{lstlisting}[style=config, caption={/etc/fail2ban/jail.local}]
[DEFAULT]
bantime  = 600    # Duree du ban : 600 secondes (10 minutes)
findtime = 300    # Fenetre d'observation : 5 minutes
maxretry = 3      # Nombre d'echecs avant ban

[samba]
enabled  = true
port     = 445,139        # Ports SMB/NetBIOS
filter   = samba          # Nom du fichier de filtre
backend  = auto
logpath  = /var/log/samba/log.*  # Tous les logs Samba
maxretry = 3
bantime  = 600
\end{lstlisting}

\textbf{Fichier} \texttt{/etc/fail2ban/filter.d/samba.conf}~:

\begin{lstlisting}[style=config, caption={/etc/fail2ban/filter.d/samba.conf}]
[Definition]
failregex = remote host \[ipv4:<HOST>:\d+\]
ignoreregex =

[Init]
datepattern = ^\[%%Y/%%m/%%d %%H:%%M:%%S
\end{lstlisting}

\begin{explainbox}[Explication détaillée de la configuration Fail2ban]
\begin{itemize}[leftmargin=1.5em,topsep=2pt]
  \item \textbf{\texttt{bantime = 600}}~: l'IP bannie est bloquée pendant 600 secondes
    (10 minutes). Après ce délai, le ban est automatiquement levé. Pour un ban permanent,
    utiliser \texttt{bantime = -1}.
  \item \textbf{\texttt{findtime = 300}}~: Fail2ban regarde en arrière sur 5 minutes. Si
    3 échecs surviennent dans cette fenêtre, le ban est déclenché.
  \item \textbf{\texttt{maxretry = 3}}~: seuil à 3 échecs avant bannissement. Choix
    délibérément bas pour stopper rapidement les attaques automatisées.
  \item \textbf{\texttt{failregex}}~: expression régulière qui identifie une ligne de log
    correspondant à un échec. \texttt{<HOST>} est remplacé par l'IP capturée.
  \item \textbf{\texttt{datepattern}}~: indique à Fail2ban le format de date dans les
    logs Samba (\texttt{[AAAA/MM/JJ HH:MM:SS]}). Les \texttt{\%\%} sont nécessaires car
    le caractère \texttt{\%} est interprété spécialement dans les fichiers ini.
  \item \textbf{Mécanisme}~: quand le seuil est atteint, Fail2ban ajoute une règle
    \texttt{nftables} (ou \texttt{iptables}) bloquant toutes les connexions depuis l'IP
    bannie sur les ports 445 et 139.
\end{itemize}
\end{explainbox}

\begin{figure}[H]
  \centering
  \includegraphics[width=0.88\textwidth]{fail2ban_marche_vraiement.png}
  \caption{Vérification du statut de Fail2ban avec \texttt{sudo systemctl status fail2ban}
    -- service \texttt{active (running)} depuis le 27 février 2026. Le PID 17831 indique
    que le démon est actif et surveille les logs en temps réel.}
\end{figure}

\subsubsection{Validation -- bannissement automatique en action}

\begin{figure}[H]
  \centering
  \includegraphics[width=\textwidth]{fail2block_kali.png}
  \caption{Fail2ban en action~: à gauche les logs montrent la progression --
    \texttt{Found 192.168.56.40} à chaque échec, puis \texttt{Ban 192.168.56.40}
    après le 3ème échec. À droite, Kali reçoit \texttt{Connection Error: The NETBIOS
    connection with the remote host timed out} -- le serveur ne répond plus du tout.}
\end{figure}

\begin{successbox}[Fail2ban opérationnel -- ban automatique confirmé]
Après exactement 3 tentatives d'authentification échouées, Fail2ban bannit
automatiquement \texttt{192.168.56.40} (Kali Linux). Les connexions suivantes de cet
attaquant sont immédiatement rejetées avec un timeout réseau, sans aucune réponse du
serveur -- ce qui ralentit aussi les scans de ports.
\end{successbox}

\begin{figure}[H]
  \centering
  \includegraphics[width=0.88\textwidth]{vraie_vue.png}
  \caption{État de la jail Samba via \texttt{sudo fail2ban-client status samba}~:
    \texttt{Total failed: 15} tentatives enregistrées, \texttt{Total banned: 1} IP
    bannie (192.168.56.40). \texttt{Currently banned: 0} confirme que le ban de 10
    minutes est expiré et que l'IP a été automatiquement libérée. L'historique reste
    consultable.}
\end{figure}

% ════════════════════════════════════════════════════════════════════════════
\section{Sauvegarde automatisée}

\subsection{Stratégie de sauvegarde}

La stratégie mise en place combine deux types de sauvegarde complémentaires~:

\begin{table}[H]
  \centering
  \begin{tabular}{lll}
    \toprule
    \textbf{Type} & \textbf{Fréquence} & \textbf{Méthode} \\
    \midrule
    Complète      & Chaque dimanche à 02h00  & \texttt{tar -czf} \\
    Incrémentielle & Lundi -- Samedi à 02h00 & \texttt{rsync --update} \\
    \bottomrule
  \end{tabular}
  \caption{Stratégie de sauvegarde automatisée}
\end{table}

\begin{explainbox}[Pourquoi combiner les deux types~?]
\begin{itemize}[leftmargin=1.5em,topsep=2pt]
  \item \textbf{Sauvegarde complète (\texttt{tar})}~: archive l'intégralité de
    \texttt{/srv/samba/} en un fichier \texttt{.tar.gz}. Point de restauration fiable
    mais lent et volumineux -- donc une fois par semaine le dimanche.
  \item \textbf{Sauvegarde incrémentielle (\texttt{rsync})}~: ne copie que les fichiers
    modifiés ou nouveaux depuis la dernière synchronisation. Rapide et économise l'espace
    disque~; exécuté tous les jours de la semaine.
  \item Cette combinaison permet de restaurer à n'importe quel point de la semaine
    (rsync quotidien) tout en ayant une archive fiable hebdomadaire (tar).
\end{itemize}
\end{explainbox}

\subsection{Script de sauvegarde}

\begin{lstlisting}[style=config, caption={/usr/local/bin/backup\_samba.sh}]
#!/bin/bash
# Script de sauvegarde automatique Samba PME

SOURCE="/srv/samba/"         # Repertoire source a sauvegarder
BACKUP_DIR="/srv/backup"     # Repertoire de destination
DATE=$(date +%Y-%m-%d)       # Format : 2026-02-27
DAY=$(date +%u)              # Jour de la semaine : 1=Lundi, 7=Dimanche

if [ "$DAY" -eq 7 ]; then
    # Dimanche : sauvegarde COMPLETE avec tar
    TYPE="complete"
    tar -czf "$BACKUP_DIR/backup_${TYPE}_${DATE}.tar.gz" "$SOURCE"
    echo "[$(date)] Sauvegarde complete effectuee" \
        >> /var/log/samba/backup.log
else
    # Lundi-Samedi : sauvegarde INCREMENTIELLE avec rsync
    TYPE="incrementielle"
    rsync -av --update "$SOURCE" "$BACKUP_DIR/incrementielle/" \
        >> /var/log/samba/backup.log 2>&1
    echo "[$(date)] Sauvegarde incrementielle effectuee" \
        >> /var/log/samba/backup.log
fi
\end{lstlisting}

\begin{explainbox}[Explication des commandes du script]
\begin{itemize}[leftmargin=1.5em,topsep=2pt]
  \item \textbf{\texttt{tar -czf fichier.tar.gz /source/}}~:
    \begin{itemize}[leftmargin=1.5em]
      \item \texttt{-c}~: crée une nouvelle archive (\textit{create})
      \item \texttt{-z}~: compresse avec \texttt{gzip} (réduit la taille d'environ 60\%)
      \item \texttt{-f}~: spécifie le nom du fichier archive
    \end{itemize}
  \item \textbf{\texttt{rsync -av --update /source/ /dest/}}~:
    \begin{itemize}[leftmargin=1.5em]
      \item \texttt{-a}~: mode archive (préserve permissions, dates, liens symboliques)
      \item \texttt{-v}~: verbeux (affiche les fichiers traités)
      \item \texttt{--update}~: ne copie que les fichiers \textit{plus récents} que leur
        correspondant en destination (incrémentiel)
    \end{itemize}
  \item \textbf{\texttt{date +\%u}}~: retourne le jour de la semaine en chiffre
    (1=lundi ... 7=dimanche). Permet de décider automatiquement le type de sauvegarde.
  \item \textbf{\texttt{>> /var/log/samba/backup.log}}~: ajoute la sortie au fichier de
    log (sans écraser les entrées précédentes).
\end{itemize}
\end{explainbox}

\subsection{Automatisation via cron}

\begin{lstlisting}[style=terminal]
# Rendre le script executable
sudo chmod +x /usr/local/bin/backup_samba.sh

# Editer la crontab de root
sudo crontab -e

# Ligne a ajouter :
# min heure jour mois jourSemaine commande
0 2 * * * /usr/local/bin/backup_samba.sh
\end{lstlisting}

\begin{figure}[H]
  \centering
  \includegraphics[width=0.88\textwidth]{cron_enregistre.png}
  \caption{Crontab root avec la tâche planifiée \texttt{0 2 * * *} -- le script
    s'exécutera tous les jours à 02h00. On voit également le contenu du journal
    \texttt{backup.log} confirmant une sauvegarde incrémentielle réussie.}
\end{figure}

\begin{explainbox}[Lecture de la syntaxe cron~: \texttt{0 2 * * *}]
La syntaxe cron comporte 5 champs~: \texttt{minute heure jourDuMois mois jourDeLaSemaine}
\begin{center}
\begin{tabular}{ccccc}
  \texttt{0} & \texttt{2} & \texttt{*} & \texttt{*} & \texttt{*} \\
  minute 0 & heure 2h & tous les jours & tous les mois & tous les jours \\
\end{tabular}
\end{center}
Le script s'exécute donc exactement à \textbf{02h00 chaque nuit}, période de faible
activité. Le script lui-même détermine ensuite si c'est un dimanche (sauvegarde
complète) ou un jour de semaine (incrémentielle).
\end{explainbox}

\subsection{Test de restauration}

La procédure de restauration a été testée en conditions réelles en quatre étapes.
Un problème de permissions a été identifié et corrigé lors de ce test, ce qui
constitue un apprentissage important.

\textbf{Étape 1 -- Création d'un fichier critique}~:
\begin{lstlisting}[style=terminal]
echo "Donnees importantes de la direction" | \
    sudo tee /srv/samba/direction/document_important.txt
\end{lstlisting}

\textbf{Étape 2 -- Sauvegarde}~:
\begin{lstlisting}[style=terminal]
sudo /usr/local/bin/backup_samba.sh
\end{lstlisting}

\begin{figure}[H]
  \centering
  \includegraphics[width=0.88\textwidth]{sauvegarde_direction.png}
  \caption{Création de \texttt{document\_important.txt} et lancement du script de
    sauvegarde. Le contenu de \texttt{/srv/backup/incrementielle/} confirme que tous
    les dossiers ont été sauvegardés.}
\end{figure}

\textbf{Étape 3 -- Simulation d'une suppression accidentelle}~:

\begin{figure}[H]
  \centering
  \includegraphics[width=0.88\textwidth]{document bien supprime.png}
  \caption{Suppression via \texttt{sudo rm} -- \texttt{document\_important.txt}
    n'apparaît plus dans le listing. Le fichier est perdu du point de vue du serveur.}
\end{figure}

% ── PROBLÈME IDENTIFIÉ ───────────────────────────────────────────────────
\subsubsection{Problème identifié~: pourquoi \texttt{sudo cp} est insuffisant}

Une première tentative de restauration avec la commande naïve \texttt{sudo cp} a été
effectuée, puis testée depuis les clients Samba~:

\begin{lstlisting}[style=terminal]
# Mauvaise pratique -- NE PAS utiliser en production
sudo cp /srv/backup/incrementielle/direction/document_important.txt \
    /srv/samba/direction/
\end{lstlisting}

\begin{figure}[H]
  \centering
  \includegraphics[width=0.88\textwidth]{permission_dossier.png}
  \caption{\texttt{sudo ls -la /srv/samba/direction/} après restauration par
    \texttt{sudo cp}~: \texttt{document\_important.txt} appartient à
    \texttt{root:direction} avec les permissions \texttt{-rw-r--r--} (644).
    Le groupe \texttt{direction} ne peut que \textbf{lire} -- le bit \texttt{w} du
    groupe est absent. Comparer avec \texttt{test\_direction} (créé directement par
    \texttt{direction1}) qui conserve les bonnes permissions avec setgid.}
\end{figure}

\begin{warningbox}[Problème critique~: \texttt{sudo cp} détruit les permissions]
\texttt{sudo cp} sans options crée le fichier restauré avec les permissions par défaut
du processus root (umask 022), donnant \texttt{-rw-r--r--} (644). Conséquences~:
\begin{itemize}[leftmargin=1.5em,topsep=2pt]
  \item \textbf{Lecture}~: \texttt{direction1} et \texttt{direction2} peuvent lire
    (bit \texttt{r} du groupe présent)~;
  \item \textbf{Écriture~: IMPOSSIBLE}~: le bit \texttt{w} du groupe est absent.
    Toute tentative de modification depuis Samba retourne
    \texttt{NT\_STATUS\_ACCESS\_DENIED}~;
  \item Le \textbf{setgid} du dossier parent (\texttt{2770}) ne corrige PAS ce
    problème~: il s'applique uniquement aux fichiers \textit{créés directement} dans
    le dossier, jamais aux fichiers \textit{copiés} par root.
\end{itemize}
\end{warningbox}

\begin{figure}[H]
  \centering
  \includegraphics[width=0.92\textwidth]{permissions_presentes.png}
  \caption{Preuve du dysfonctionnement~: \texttt{direction1} \textbf{et} \texttt{direction2}
    obtiennent \texttt{NT\_STATUS\_ACCESS\_DENIED opening remote file
    \textbackslash document\_important.txt} en tentant d'écrire via smbclient.
    Aucun membre du service Direction ne peut modifier le fichier restauré.
    Note~: \texttt{rh1} reçoit \texttt{tree connect failed} (normal -- il
    n'appartient pas au groupe \texttt{direction} et ne peut pas accéder au
    partage). La restauration par \texttt{sudo cp} est donc inutilisable en
    conditions réelles.}
\end{figure}

% ── SOLUTION CORRECTE ────────────────────────────────────────────────────
\subsubsection{Étape 4 -- Restauration correcte avec \texttt{rsync --archive}}

La solution est d'utiliser \texttt{rsync} avec le flag \texttt{--archive}, qui
\textbf{préserve les permissions, le groupe propriétaire et les dates} tels qu'ils
étaient au moment de la sauvegarde~:

\begin{lstlisting}[style=terminal]
# Bonne pratique -- preserve les permissions et le groupe
sudo rsync -av --archive \
    /srv/backup/incrementielle/direction/document_important.txt \
    /srv/samba/direction/
\end{lstlisting}

\begin{figure}[H]
  \centering
  \includegraphics[width=0.88\textwidth]{restauration_correcte.png}
  \caption{\texttt{rsync -av --archive} en action~: \texttt{sending incremental
    file list} confirme l'analyse des différences, puis \texttt{document\_important.txt}
    indique le fichier transféré. Les statistiques (\texttt{sent 78 bytes, received
    12 bytes, total size is 37}) confirment une restauration réussie avec
    préservation complète des métadonnées.}
\end{figure}

\begin{explainbox}[Explication détaillée~: pourquoi \texttt{rsync -av --archive} est la bonne commande]
\begin{itemize}[leftmargin=1.5em,topsep=2pt]
  \item \texttt{rsync}~: outil de synchronisation qui compare source et destination
    avant de copier, ne transférant que ce qui a changé.
  \item \texttt{--archive} (ou \texttt{-a})~: mode archive -- combine en une seule
    option~:
    \begin{itemize}[leftmargin=1.5em]
      \item \texttt{-r}~: récursif (sous-dossiers inclus)
      \item \texttt{-l}~: préserve les liens symboliques
      \item \texttt{-p}~: préserve les \textbf{permissions} (chmod) -- résout le
        problème des 644 vs 660
      \item \texttt{-t}~: préserve les \textbf{dates} de modification
      \item \texttt{-g}~: préserve le \textbf{groupe propriétaire} -- le fichier
        retrouve son groupe \texttt{direction} au lieu de \texttt{root}
      \item \texttt{-o}~: préserve le propriétaire (nécessite root)
    \end{itemize}
  \item \texttt{-v}~: mode verbeux, affiche les fichiers traités.
  \item \textbf{Différence fondamentale avec \texttt{cp}}~: \texttt{cp} sans
    \texttt{-a} ou \texttt{-p} crée un nouveau fichier avec les permissions par
    défaut du processus appelant (root $\rightarrow$ 644 avec umask 022).
    \texttt{rsync --archive} reproduit le fichier \textit{à l'identique}, permissions
    et groupe propriétaire inclus.
\end{itemize}
\end{explainbox}

\begin{successbox}[Restauration correcte validée]
Après \texttt{rsync --archive}, le fichier restauré appartient à \texttt{root:direction}
avec les permissions \texttt{-rw-rw----} (660)~: les membres du groupe
\texttt{direction} peuvent lire \textbf{et écrire}. La restauration est complète et
opérationnelle pour tous les utilisateurs du service.\\[6pt]
\textbf{Pour restaurer un dossier entier}~:\\
\texttt{sudo rsync -av --archive /srv/backup/incrementielle/direction/}
\texttt{/srv/samba/direction/}
\end{successbox}


% ════════════════════════════════════════════════════════════════════════════
\section{Difficultés rencontrées et choix techniques}

\subsection{Difficultés rencontrées}

\subsubsection{Configuration réseau de la VM}

La principale difficulté initiale a été la configuration du réseau. L'interface Host-Only
(\texttt{enp0s8}) restait en état \texttt{NO-CARRIER} même après la configuration netplan.
La cause~: l'option \textit{Câble connecté} n'était pas cochée dans les paramètres VirtualBox
de l'adaptateur réseau. Une fois cette option activée et l'IP configurée en statique
(\texttt{192.168.56.20/24}), l'interface est devenue fonctionnelle.

\subsubsection{Instabilité de l'IP de Kali}

Les IPs assignées manuellement avec \texttt{ip addr add} sur Kali ne sont pas persistantes
-- elles disparaissent à chaque redémarrage ou coupure de courant (qui est survenue plusieurs
fois lors du travail). Cela a nécessité de réassigner l'IP plusieurs fois pendant les tests,
et a obligé à changer l'adresse de \texttt{192.168.56.30} à \texttt{192.168.56.40} à un
moment après un conflit. La solution pérenne aurait été de configurer la persistance via
\texttt{/etc/network/interfaces} sur Kali.

\subsubsection{Configuration de Fail2ban}

Le filtre Fail2ban a nécessité plusieurs itérations. Les problèmes rencontrés~:
\begin{enumerate}[leftmargin=2em]
  \item Le caractère \texttt{\%} dans les patterns de date est interprété différemment par
    Fail2ban (nécessite \texttt{\%\%} dans la section \texttt{[Init]})~;
  \item Le filtre initial ne contenait pas de groupe \texttt{<HOST>}, rendant impossible
    l'identification de l'IP à bannir~;
  \item Les logs Samba n'incluent pas l'IP dans les lignes d'échec direct -- elle se trouve
    dans les lignes \texttt{Auth:} avec le format \texttt{remote host [ipv4:X.X.X.X:port]}~;
  \item La directive \texttt{logpath = /var/log/samba/log.\%m} contenait \texttt{\%m} que
    Fail2ban interprète comme une variable -- remplacement par \texttt{log.*} (glob) nécessaire.
\end{enumerate}

\subsubsection{Conflit d'IP avec le réseau Host-Only}

Lors de la première configuration, l'IP \texttt{192.168.56.1} avait été assignée à la VM
alors que c'est l'IP par défaut de l'interface \texttt{vboxnet0} du PC hôte. Le conflit
générait des erreurs de routage. Résolu en assignant \texttt{192.168.56.20} à la VM.

\subsection{Choix techniques justifiés}

\begin{table}[H]
  \centering
  \begin{tabular}{p{3.8cm}p{3.8cm}p{5.5cm}}
    \toprule
    \textbf{Choix} & \textbf{Alternative} & \textbf{Justification} \\
    \midrule
    Ubuntu Server 24.04 & Debian 12 & Déjà disponible et fonctionnel, moins de configuration initiale, meilleures docs en ligne \\
    Host-Only + Bridge  & NAT seul  & Bridge pour Internet (\texttt{apt}), Host-Only pour isoler les tests d'attaque \\
    \texttt{rsync} (incrémentiel) & \texttt{tar} complet quotidien & Plus rapide, ne sauvegarde que les fichiers modifiés, économise l'espace \\
    \texttt{fail2ban} & \texttt{ufw} rate-limit & Fail2ban analyse les logs applicatifs (Samba) et non juste les paquets réseau \\
    \texttt{nftables} (via fail2ban) & \texttt{iptables} & Standard moderne sur Ubuntu 24, meilleure performance et syntaxe unifiée \\
    Mot de passe complexe & Mot de passe simple & Résiste aux attaques par dictionnaire comme \texttt{rockyou.txt} \\
    \bottomrule
  \end{tabular}
  \caption{Justification des choix techniques}
\end{table}

% ════════════════════════════════════════════════════════════════════════════
\section{Perspectives d'amélioration}

Plusieurs améliorations pourraient renforcer ce déploiement en conditions de production~:

\begin{enumerate}[leftmargin=2em]
  \item \textbf{Chiffrement des communications}~: activer le chiffrement SMB3
    (\texttt{smb encrypt = required}) pour chiffrer toutes les données en transit.
    Empêche l'interception même sur un réseau compromis~;

  \item \textbf{Intégration LDAP/Active Directory}~: remplacer la gestion locale des
    utilisateurs par un annuaire centralisé. Indispensable au-delà d'une dizaine
    d'utilisateurs pour éviter la gestion manuelle~;

  \item \textbf{Restriction d'accès par IP}~: ajouter \texttt{hosts allow = 192.168.56.}
    dans \texttt{smb.conf} pour limiter les connexions aux seules machines du réseau interne~;

  \item \textbf{Sauvegardes distantes}~: envoyer les sauvegardes vers un serveur distant
    ou un stockage cloud (S3, NFS) pour garantir la disponibilité en cas de sinistre
    physique (incendie, vol) sur le site~;

  \item \textbf{Persistance des IPs sur Kali}~: configurer \texttt{/etc/network/interfaces}
    pour rendre les IPs persistantes entre les redémarrages, évitant de reconfigurer
    à chaque session de test~;

  \item \textbf{Bannissement permanent des récidivistes}~: passer \texttt{bantime = -1}
    dans Fail2ban avec une whitelist des IPs légit pour les attaques répétées~;

  \item \textbf{Tests de sauvegarde complète}~: bien que la sauvegarde incrémentielle ait
    été testée et validée, il serait important de valider également la restauration depuis
    une archive \texttt{tar} complète (test de disaster recovery).
\end{enumerate}

% ════════════════════════════════════════════════════════════════════════════
\section{Conclusion}

Ce projet a permis de mettre en place un serveur de fichiers Samba complet et sécurisé
pour une PME fictive, en partant de zéro jusqu'à un déploiement opérationnel et testé.
L'ensemble des exigences du cahier des charges a été satisfait~:

\begin{itemize}[leftmargin=2em]
  \item Quatre groupes et huit utilisateurs créés avec des espaces dédiés et isolés~;
  \item Permissions strictement différenciées par groupe, avec espace commun partagé~;
  \item SMBv1 désactivé, communications en SMB2/SMB3 uniquement~;
  \item Journalisation complète avec traçabilité de toutes les connexions~;
  \item Protection contre l'énumération anonyme (\texttt{restrict anonymous = 2})~;
  \item Résistance démontrée aux attaques par brute force (Hydra échoue, CrackMapExec bloqué)~;
  \item Bannissement automatique des attaquants via Fail2ban (validé en conditions réelles)~;
  \item Stratégie de sauvegarde complète + incrémentielle automatisée via cron~;
  \item Procédure de restauration testée et validée en conditions réelles.
\end{itemize}

Les difficultés rencontrées -- notamment la configuration de Fail2ban (plusieurs itérations
pour obtenir le bon filtre regex) et les instabilités réseau liées aux coupures de courant
-- ont été autant d'occasions d'approfondir la compréhension des mécanismes sous-jacents. Ce déploiement, réalisé en environnement virtualisé, constitue une
base solide transposable à un environnement de production réel.

% ════════════════════════════════════════════════════════════════════════════
\section*{Références}

\begin{enumerate}[leftmargin=2em]
  \item Samba Team. \textit{Samba -- Opening Windows to a Wider World}.
    \url{https://www.samba.org/}
  \item Microsoft. \textit{SMB security enhancements}.
    \url{https://learn.microsoft.com/en-us/windows-server/storage/file-server/smb-security}
  \item Fail2ban Project. \url{https://www.fail2ban.org/}
  \item Ubuntu Server Documentation. \url{https://ubuntu.com/server/docs}
  \item Netplan Reference. \url{https://netplan.io/reference/}
  \item The Linux Documentation Project. \textit{Linux File Permissions}.
    \url{https://tldp.org/LDP/intro-linux/html/sect_03_04.html}
\end{enumerate}

% ════════════════════════════════════════════════════════════════════════════
\appendix
\section{Annexe -- Récapitulatif des commandes utilisées}

\begin{longtable}{ll}
  \toprule
  \textbf{Commande} & \textbf{Description} \\
  \midrule
  \endhead
  \texttt{ip a} & Afficher les interfaces et adresses IP \\
  \texttt{sudo netplan apply} & Appliquer la configuration réseau \\
  \texttt{sudo apt install samba} & Installer Samba \\
  \texttt{samba --version} & Vérifier la version de Samba \\
  \texttt{sudo systemctl status smbd} & Statut du service Samba \\
  \texttt{sudo systemctl restart smbd} & Redémarrer le service Samba \\
  \texttt{sudo groupadd <groupe>} & Créer un groupe Linux \\
  \texttt{sudo useradd -m -G <groupe> <user>} & Créer un utilisateur dans un groupe \\
  \texttt{sudo passwd <user>} & Définir le mot de passe Linux \\
  \texttt{sudo smbpasswd -a <user>} & Enregistrer l'utilisateur dans Samba \\
  \texttt{sudo chown root:<groupe> <dossier>} & Changer le groupe propriétaire \\
  \texttt{sudo chmod 2770 <dossier>} & Permissions setgid + accès total groupe \\
  \texttt{sudo chmod 2777 <dossier>} & Permissions setgid + accès total tous \\
  \texttt{testparm} & Valider la syntaxe de smb.conf \\
  \texttt{smbclient -L //IP -N} & Lister les partages en anonyme (test) \\
  \texttt{smbclient -L //IP -U user} & Lister les partages avec authentification \\
  \texttt{smbclient //IP/share -U user} & Se connecter à un partage \\
  \texttt{sudo tail -f /var/log/samba/log.*} & Surveiller les logs en temps réel \\
  \texttt{sudo fail2ban-client status samba} & Statut de la jail Samba dans Fail2ban \\
  \texttt{sudo fail2ban-regex <log> <filter>} & Tester un filtre Fail2ban \\
  \texttt{sudo /usr/local/bin/backup\_samba.sh} & Lancer manuellement la sauvegarde \\
  \texttt{sudo crontab -l} & Lister les tâches cron root \\
  \texttt{sudo crontab -e} & Éditer les tâches cron root \\
  \texttt{tar -czf archive.tar.gz /dossier/} & Créer une archive compressée \\
  \texttt{rsync -av --update /src/ /dst/} & Synchronisation incrémentielle \\
  \bottomrule
\end{longtable}

\end{document}